% @Author: AnthonyKenny98
% @Date:   2019-10-30 22:08:06
% @Last Modified by:   AnthonyKenny98
% @Last Modified time: 2020-04-03 10:09:41

% Define Document Class
\documentclass[
    11pt,           % Font Size
    letterpaper,    % Page Size
    % draft,          % Draft Mode - error checking and no image rendering
    oneside         % Single Sided
]{report}           % Report Mode

% Import custom commands
\RequirePackage{../support/thesis}
\RequirePackage{../support/codestyle}

\usepackage[acronym]{glossaries}
\makeglossaries

\newcommand{\definenewglossaryterm}[4]{
    % #1: Identifier and Acronym e.g. API
    % #2: Long e.g. Application Programming Interface
    % #3: Definition. e.g. A particular set of rules and specifications...


    % \newglossaryentry{#1}{
    %     text={#2},
    %     long={#3},
    %     name={\glsentrylong{#1} (\glsentrytext{#1})},
    %     first={\glsentryname{#1}},
    %     firstplural={\glsentrylong{#1}\glspluralsuffix (\glsentryname{#1}\glspluralsuffix )},
    %     description={#4}
    % }

    \newglossaryentry{#1_g}{
        text={#1},
        long={#2},
        name={\glsentrylong{#1_g} (\glsentrytext{#1}):},
        description={#3}}

    \newglossaryentry{#1}{
        type=\acronymtype,
        name={#1},
        description={#2},
        first={#2 (#1)\glsadd{#1_g}}, see=[Glossary:]{#1_g},
        firstplural={\glsentrylong{#1}\glspluralsuffix (\glsentryname{#1}\glspluralsuffix)},
        plural={\glsentryname{#1}\glspluralsuffix}
    }
}

\definenewglossaryterm{API}{Application Programming Interface}{A particular set of rules and specifications that a software program can follow to access and make use of the services and resources provided by another particular software program that implements that API}

% \newglossaryentry{apig}{name={API},
%     description={An Application Programming Interface (API) is a particular set
% of rules and specifications that a software program can follow to access and
% make use of the services and resources provided by another particular software
% program that implements that API}}

% %%% define the acronym and use the see= option
% \newglossaryentry{api}{type=\acronymtype, name={API}, description={Application
% Programming Interface}, first={Application
% Programming Interface (API)\glsadd{apig}}, see=[Glossary:]{apig}}

%%%%%%%%%%%%%%%%%%%%%%%%%%%%%%%%%%%%%%%%%%%%%%%%%%%%%%%%%%

% DEFINE HEADER & FOOTER: (Name and Page Number)
\pagestyle{fancy}
\setlength{\headheight}{14pt}
\fancyhf{}
\rhead{Anthony J.W. Kenny}
\cfoot{\thepage}

% DEFINE COVER PAGE
\title{\textsc{if goingToCrash(): dont()\ \\ RISC-V Architecture for Motion Planning Algorithms in Autonomous UAVs} \\ 
    \bigskip
    \small{A senior design project submitted in partial fulfillment of the requirements for the degree of Bachelor of Science at Harvard University} \\
\author{Anthony J.W. Kenny \\
        \small{S.B. Candidate in Electrical Engineering} \\ \\
        Faculty Advisor: Vijay Janapa Reddi \\ \\
        Harvard University School of Engineering and Applied Sciences \\
        \small{Cambridge, MA}}
\date{March 1st, 2020 \\ 
    \small{Version: 4.2}}}

\begin{document}

% Cover Page
\maketitle

% PREFACE
\pagenumbering{roman}
\addcontentsline{toc}{chapter}{Preface}

    % ABSTRACT
    \addcontentsline{toc}{section}{Abstract}
    % @Author: AnthonyKenny98
% @Date:   2020-04-03 15:24:10
% @Last Modified by:   AnthonyKenny98
% @Last Modified time: 2020-04-10 11:33:34
\begin{abstract}
\doublespacing
This thesis presents an extension for the RISC-V Instruction Set Architecture (ISA) for the purpose of faster motion planning in autonomous Unmanned Aerial Vehicles (UAVs). Fully autonomous UAVs have the potential to change the world in which we live, but they are currently unable to pilot themselves in high-complexity, obstacle-dense environments; The processors that they employ cannot execute motion planning software quickly enough. RISC-V is a relatively new ISA that is founded on the principles of open-source and extendibility, making it an excellent ecosystem for designing application specific processors. However, as of April 2020, no attempts had been made to develop motion planning architecture within the RISC-V ecosystem. \\
This thesis serves as a proof-of-concept for accelerating motion planning with RISC-V architecture. It presents the implementation of microarchitecture that can compute collision detection over 3 times faster than a general-purpose Intel CPU. More significantly, it defines a motion planning extension for RISC-V that simplifies the number of instructions required to detect and edge collision from hundreds of thousands to only one. These promising results demonstrate the viability of the approach, and should encourage developers to embrace RISC-V in the development of motion planning processors for autonomous drones.

\end{abstract}
    \todo[inline]{Rewrite Abstract}
    \clearpage

    % TABLE OF CONTENTS
    \addcontentsline{toc}{section}{Table of Contents}
    \tableofcontents
    \clearpage

    % ACRONYMNS
    \section*{List of Acronyms}
    \addcontentsline{toc}{section}{List of Acronyms}
    % \input{\acronymsName}

    % ALGORITHMS
    % \listofalgorithms
 
    % FIGURES
    % \listoffigures

    % TABLES
    % \listoftables

    \newpage
    First Use \gls{API}\\ 
    Second Use \gls{API} 
    Plural Use \glspl{API}

    \newpage
    \printglossary[type=\acronymtype]
    \newpage
    \printglossary[type=main]


% CHAPTERS

% Chapter 1
\chapter{Introduction}
    \pagenumbering{arabic}                      % Change Page numbering to arabic
    % @Author: AnthonyKenny98
% @Date:   2020-02-20 00:00:23
% @Last Modified by:   AnthonyKenny98
% @Last Modified time: 2020-02-20 01:19:31


This is a pretend citation\cite{Murrayb} and this is an incorrectly spelled words.

\newpage
Now my chapter 1 is two pages
    \lhead{Chapter \thechapter}                 % Include Chapter name on subsequent headers


% Chapter 2
\chapter{Motion Planning in Software}
    \label{chap:MotionPlanningInSoftware}
    % @Author: AnthonyKenny98
% @Date:   2020-02-22 15:42:12
% @Last Modified by:   AnthonyKenny98
% @Last Modified time: 2020-04-05 17:17:11

The first objective of this thesis is to identify a typical motion planning algorithm, profile its execution, and determine computational bottlenecks.\todo{Update once I have properly defined goals and objectives}

This chapter introduces the concept of motion planning and details the process of implementing and analysing \glsfirst{RRT}, a commonly used algorithm, to identify its computational bottlenecks.\\

\section{Motion Planning Background} 
\label{section:motion_planning_background}
    % @Author: AnthonyKenny98
% @Date:   2020-04-04 10:01:40
% @Last Modified by:   AnthonyKenny98
% @Last Modified time: 2020-04-04 11:17:14

A funny paradox in computer science is the fact that it is relatively easy to teach a computer to perform tasks that humans find very complicated, but extremely difficult to program one to execute functions that humans master during infancy. Consider, it was as early as 1949 that Claude Shannon presented his paper \textit{Programming a Computer for Playing Chess}\cite{Shannon1950}, and by 1997 the \textit{Deep Blue} computer defeated Garry Kasparov, reigning world champion, in a six game chess match.\cite{Campbell2002} Compare that with some of the most advanced autonomous humanoid robots to date displaying dexterity comparable with a toddler. The task of finding a collision free path, performed constantly without thought by a human, is an example of this paradigm. For a robot to compute a collision free path, it relies on a set of Motion Planning Algorithms.

Motion Planning Algorithms refer to the set of algorithms that find possible sequences of valid \gls{configuration}s for a robot in a space. In plain English, they are algorithms that determine the movements a robot can make in a map, with the intent of eventually finding a path from one point to another. 

\subsection{Key Concepts}

    \subsubsection{Configuration}

    % @Author: AnthonyKenny98
% @Date:   2020-04-04 11:13:58
% @Last Modified by:   AnthonyKenny98
% @Last Modified time: 2020-04-04 11:23:02
% @Author: AnthonyKenny98
% @Date:   2020-02-29 17:30:44
% @Last Modified by:   AnthonyKenny98
% @Last Modified time: 2020-04-03 14:28:30
\begin{figure}[H]
\begin{center}
\begin{tabular}{cc}

    % Subfigure A
    \begin{subfigure}{\textwidth}
    \begin{center}
    \includegraphics[width=0.4\linewidth]{chapters/chapter2/img/motionPlanning/3DPointConfiguration.png}
    \caption{A robot represented by just a point in 3D space, requiring only 3 Cartesian coordinate $(x,y,z)$ points to describe its \gls{configuration}}
    \label{subfig:3DPointConfig}
    \end{center}
    \end{subfigure}
    &
    % 
    % Subfigure B
    \begin{subfigure}{\textwidth}
    \begin{center}
    \includegraphics[width=0.4\linewidth]{chapters/chapter2/img/motionPlanning/3DCubeConfiguration.png}
    \caption{A robot represented as a cube in 3D space, now requiring 3 Euler angles $(\alpha, \beta, \gamma)$ along with the original Cartesian coordinates.}
    \label{subfig:3DCubeConfig}
    \end{center}
    \end{subfigure} \\
\end{tabular}
    % Caption and Label
    \caption{Configuration TODO}
    \label{fig:configuration}

\end{center}
\end{figure}



\newpage
\section{Implementation of RRT}
\label{section:rrt}
    % @Author: AnthonyKenny98
% @Date:   2020-02-22 15:53:59
% @Last Modified by:   AnthonyKenny98
% @Last Modified time: 2020-04-05 18:38:15


% TECHNICAL SPECIFICATIONS
\subsection{Technical Specifications}

    With \gls{RRT} selected as the benchmark algorithm against which to test specialised hardware, this project required an implementation of the algorithm that satisfied the following criteria shown in Table \ref{table:RRT_Tech_Specs_Abbrev}. Appendix \ref{section:rrt_appendix_tech_specs} is a more thorough description of the technical specifications for the implementation of RRT. 

    % Tech Spec Table
    \begin{table}[H]
\begin{center}
\begin{tabular}{|p{.3\linewidth}|p{.64\linewidth}|}
    \hline
    Requirement             & Description and Justification \\
    \hline
    C/C++ Implementation    & As outlined in Section \ref{subsection:project_structure}, the critical step in determining the design of specialized hardware to accelerate \gls{RRT} is CPU performance analysis of the algorithm to determine computational hot-spots. Implementations in C allow for the use of certain CPU profiling tools, described in Section \ref{subsubsection:vtune}, unlike higher-level languages such as Python. \\
    \hline
    Models Drone as Point   & In reality, implementing \gls{RRT} for a drone would model the robot as a \gls{3D} object defined by coordinates $\{x, y, z\}$ and Euler angles $\{\alpha, \beta, \gamma \}$. However, for simplicity's sake, modelling the drone as a point defined by coordinates $\{x, y, z\}$ will suffice. Time permitting, this could be revisited. \todo[inline]{Change this based on whether time does permit} \\
    \hline
    Mirrors Algorithm       & In order for the results of CPU performance analysis to be easy to understand, software implementation of \gls{RRT} should call functions that mirror the functions described in Algorithms \ref{algorithm:rrt} and \ref{algorithm:rrt_collision}. \\
    \hline
\end{tabular}
\caption{Technical Specifications for \gls{RRT} Implementation}
\label{table:RRT_Tech_Specs}
\end{center}
\end{table}
\todo[inline]{Improve this table}

    The original intention was to find an existing implementation of RRT that could fulfill these requirements. However, no open-source implementations were suitable. Appendix \ref{section:rrt_appendix_existing_implementations} shows an evaluation of existing implementations.

    As a result, it was necessary to build a C implementation of RRT from the ground up to the aformentioned specifications.

\subsection{Implementation Design}
    
    The design and implementation of \gls{RRT}, while neccessary, was significant and time consuming. Since this was not the main object of this thesis, only a brief description of key design choices has been included here. Appendix \ref{appendix:rrt_appendix} contains a more detailed account.

    \subsubsection{Parameterization}
        Table \ref{table:rrt_params} shows the parameters that were included in the implementation and compiled by way of a C header file.
        % @Author: AnthonyKenny98
% @Date:   2020-04-05 17:54:44
% @Last Modified by:   AnthonyKenny98
% @Last Modified time: 2020-04-05 18:10:00
\begin{table}[H]
\begin{centering}
\begin{tabular}{|c|c|l|}
\hline
\textbf{Parameter} & \textbf{Data Type} & \textbf{Description} \\
\hline
$\epsilon$ & Integer & Maximum distance between two configurations \\
\hline
$K$ & Integer & Maximum number of configurations in the graph \\
\hline
$DIM$ & Integer & Upper bound of each axis of workspace \\
\hline
Goal Bias & Float & Percentage likelihood of stepping towards goal node \\
\hline
\gls{OGM} & File Pointer & \glsname{CSV} of booleans to represent grids \\
\hline
\end{tabular}
\caption{RRT Implementation Parameters}
\label{table:rrt_params}

\end{centering}
\end{table}

    \subsubsection{Key Functions}
        Algorithm \ref{algorithm:rrt_collision} shows that there are 5 key functions that constitute \gls{RRT}. Figure \ref{fig:rrt_functions} demonstrates each of these functions: \texttt{getRandomConfig()}, \texttt{findNearestConfig()}, \texttt{stepFromNearest()}, \texttt{(configCollisions)}, and \texttt{edgeCollisions()}.
        % @Author: AnthonyKenny98
% @Date:   2020-04-05 18:33:38
% @Last Modified by:   AnthonyKenny98
% @Last Modified time: 2020-04-05 18:43:04
\begin{figure}[H]
\begin{centering}
\begin{tabular}{ccc}

    \begin{subfigure}{0.3\linewidth}
    \includegraphics[width=\linewidth]{chapters/chapter2/img/keyfunctions/functions1.png}
    \caption{}
    \end{subfigure} & 

    \begin{subfigure}{0.3\linewidth}
    \includegraphics[width=\linewidth]{chapters/chapter2/img/keyfunctions/functions2.png}
    \caption{}
    \end{subfigure} &

    \begin{subfigure}{0.3\linewidth}
    \includegraphics[width=\linewidth]{chapters/chapter2/img/keyfunctions/functions3.png}
    \caption{}
    \end{subfigure} \\

    \begin{subfigure}{0.3\linewidth}
    \includegraphics[width=\linewidth]{chapters/chapter2/img/keyfunctions/functions4.png}
    \caption{}
    \end{subfigure} &

    \begin{subfigure}{0.3\linewidth}
    \includegraphics[width=\linewidth]{chapters/chapter2/img/keyfunctions/functions5.png}
    \caption{}
    \end{subfigure} & 

    \begin{subfigure}{0.3\linewidth}
    \includegraphics[width=\linewidth]{chapters/chapter2/img/keyfunctions/functions6.png}
    \caption{}
    \end{subfigure} \\

\end{tabular}
\caption{}
\label{fig:rrt_functions}
\end{centering}
\end{figure}

    \subsubsection{Dimensionality}

    \subsubsection{$K$:DIM Ratio}

\subsection{Visualizing Implementation}
% @Author: AnthonyKenny98
% @Date:   2020-02-23 14:14:12
% @Last Modified by:   AnthonyKenny98
% @Last Modified time: 2020-03-01 08:06:12


\begin{figure}[H]
\begin{center}
\begin{tabular}{c  c}
    \includegraphics[width=0.45\linewidth]{chapters/chapter2/img/rrt_2d_1.png} & \includegraphics[width=0.45\linewidth]{chapters/chapter2/img/rrt_2d_2.png} \\
    \includegraphics[width=0.45\linewidth]{chapters/chapter2/img/rrt_2d_3.png} & \includegraphics[width=0.45\linewidth]{chapters/chapter2/img/rrt_2d_4.png}
    \end{tabular}
    \caption{2D RRT Implementation shown by \ac{GUI}}
    \label{figure:2DrrtGui}
\end{center}
\end{figure}

\todo[inline]{Describe implementation in 3D}
% @Author: AnthonyKenny98
% @Date:   2020-02-23 14:14:12
% @Last Modified by:   AnthonyKenny98
% @Last Modified time: 2020-02-23 14:19:03


\begin{figure}[H]
\begin{center}
    % \begin{tabular}{c  c}
    % \includegraphics[draft=false,width=0.45\linewidth]{chapters/chapter2/img/rrt_2d_1.png} & \includegraphics[draft=false,width=0.45\linewidth]{chapters/chapter2/img/rrt_2d_2.png} \\
    % \includegraphics[draft=false,width=0.45\linewidth]{chapters/chapter2/img/rrt_2d_3.png} & \includegraphics[draft=false,width=0.45\linewidth]{chapters/chapter2/img/rrt_2d_4.png}
    % \end{tabular}
    \missingfigure[figwidth=\linewidth, figheight=10cm]{Simulations of 3D RRT Implementations}
    \caption{3D RRT Implementation shown by \ac{GUI}}
    \label{figure:3DrrtGui}
\end{center}
\end{figure}


\newpage
\section{Analysis of RRT}
\label{section:rrt_analysis}
    % @Author: AnthonyKenny98
% @Date:   2020-02-28 15:02:19
% @Last Modified by:   AnthonyKenny98
% @Last Modified time: 2020-03-01 15:49:31

\todo[inline]{Brief introduction outlining purpose of performance analysis}

\subsection{Methodology}
    To restate, the aim of this thesis is to design a computer processor with reduced execution time of motion planning algorithms, such as \ac{RRT}. As such, it is important to understand the elements of the algorithm that have the highest percentage of CPU execution time. To determine this, it was necessary to implement my own, naive but typical, \ac{RRT} in C. This program could then be compiled and analysed using a software performance profiling tool. With this, I could design experiments to determine the critical RRT functions (those occupying a majority of CPU time) and see how this varies given different parameters.
    \todo[inline]{Outline of method of analysis. Something better than the above}

    \subsubsection*{VTune Profiler}
    \label{subsubsection:vtune}
        VTune Profiler performance profiler is an application for software performance analysis. It provides functionality to examine hot-spots for CPU execution time through a top down analysis, shown below in Figure \ref{figure:VTuneTopDown}. As can be seen from the figure, the top down analysis tool shows the percentage of CPU time taken up by each function. I used this tool to profile the algorithm's performance as I changed certain parameters.
        \todo[inline]{Rewrite the above}
        % @Author: AnthonyKenny98
% @Date:   2020-02-23 14:33:19
% @Last Modified by:   AnthonyKenny98
% @Last Modified time: 2020-02-23 14:36:06

\begin{figure}[H]
\begin{center}
    \missingfigure[figwidth=\linewidth]{Screenshot of VTune Top Down Analysis (Maybe)}
    \caption{VTune Amplifier TopDown Analysis Example}
    \label{figure:VTuneTopDown}
\end{center}
\end{figure}

    \subsubsection*{Internal Timing}
        The limitation of VTune Profiler is that it can only profile software running on Intel processors, which implement the x86-64 \ac{ISA}. As such, when the time comes to analyse performance of the software running on a RISC-V processor, another method will be required. A simple and effective way of measuring execution performance is to insert timing functionality into the software itself. \\

        \todo[inline]{Provide or link to appendix of explanation of internal timing}

    \subsubsection*{Comparison}
        Before proceeding to use either of these methods to profile the software implementation of \ac{RRT}, it was important to verify that the two methods yielded similar results for the same program. Table \ref{table:timing_calibration} summarizes the results of analysis of a simple C executable. The program calls 5 functions, $\{A, B, C, D, E\}$, each a simple iteration in which a integer is incremented. Since the Internal Timing method returned similar results to the (trusted) VTune Profiler, it was considered to be a reliable method. While it was encouraging to see both methods returned similar results for absolute execution time, the more important metric was the similarity in percentage of total execution time.

        \begin{table}[H]
\begin{center}
\begin{tabular}{|m{0.1\linewidth}|m{0.18\linewidth}|m{0.18\linewidth}|m{0.18\linewidth}|m{0.18\linewidth}|}
\hline
\multirow{2}{*}{function} & \multicolumn{2}{c|}{Vtune Profiler} & \multicolumn{2}{c|}{Internal Timing} \\
\cline{2-5}
            & time (s)  & time (\% total) & time (s) & time (\% total)    \\
\hline
A       & 0.488     & 57.4\%        & 0.497 & 57.6\%           \\
B       & 0.2       & 23.5\%        & 0.198 & 23.1\%           \\
C       & 0.102     & 12.0\%        & 0.099 & 11.5\%           \\
D       & 0.048     & 5.7\%         & 0.049 & 5.6\%            \\
E       & 0.012     & 1.4\%         & 0.019 & 2.2\%            \\
\hline
\end{tabular}
\end{center}
\caption{Comparison of Timing Methods}
\label{table:timing_calibration}
\end{table}

    \subsubsection*{Experimental Design}
        In profiling \ac{RRT} in software, the goal was to find the critical task across different values of $K$ and sizes of configuration space. Multiple tests were run, varying these two constraints, to find this critical function. The results of this analysis can be found in Section \ref{section:rrt_analysis_results}.

\subsection{Results}
\label{section:rrt_analysis_results}
    Figure \ref{fig:rrt_profiling} shows the profile of functions within \ac{RRT}, for $100 \leq K \leq 10000$, and cubic configuration spaces with dimensions $\{4, 8, 16, 32\}$. Each subfigure shows a similar profile, with the \% of CPU Execution Time taken by findNearestNode increasing with $K$. This is to be expected. \todo{Explanation of time complexity}However, it is also seen that edgeCollisions increases with larger configuration spaces, taking up the overwhelming majority of execution time for a 32x32x32 configuration space.
    
    \newpage
    % @Author: AnthonyKenny98
% @Date:   2020-02-29 14:51:44
% @Last Modified by:   AnthonyKenny98
% @Last Modified time: 2020-04-03 14:28:30
\begin{figure}[H]
\begin{center}

    % Subfigure A
    \begin{subfigure}{\textwidth}
    \begin{center}
    \includegraphics[width=\linewidth,height=0.3\paperheight]{chapters/chapter2/img/profiling/4x4x4/performance.png}
    \caption{4x4x4 \gls{configuration} Space}
    \label{subfig:4x4x4rrt}
    \end{center}
    \end{subfigure}
    % 
    % Subfigure B
    \begin{subfigure}{\textwidth}
    \begin{center}
    \includegraphics[width=\linewidth,height=0.3\paperheight]{chapters/chapter2/img/profiling/8x8x8/performance.png}
    \caption{8x8x8 \gls{configuration} Space}
    \label{subfig:8x8x8rrt}
    \end{center}
    \end{subfigure}
    % Caption and Label
    \caption{\gls{RRT} Functions as a \% of Total CPU Exectution Time}
\end{center}
\end{figure}

\newpage
\begin{figure}[H]\ContinuedFloat
\begin{center}
    % 
    % Subfigure C
    \begin{subfigure}{\textwidth}
    \begin{center}
    \includegraphics[width=\linewidth,height=0.3\paperheight]{chapters/chapter2/img/profiling/16x16x16/performance.png}
    \caption{16x16x16 \gls{configuration} Space}
    \label{subfig:16x16x16rrt}
    \end{center}
    \end{subfigure}
    % 
    % Subfigure D
    \begin{subfigure}{\textwidth}
    \begin{center}
    \includegraphics[width=\linewidth,height=0.3\paperheight]{chapters/chapter2/img/profiling/32x32x32/performance.png}
    \caption{32x32x32 \gls{configuration} Space}
    \label{subfig:32x32x32rrt}
    \end{center}
    \end{subfigure} 
    
    % Caption and Label
    \caption{\gls{RRT} Functions as a \% of Total CPU Exectution Time (cont.)}
    \label{fig:rrt_profiling}
\end{center}
\end{figure}
\todo[inline]{Change Y axis to \% and increase text size}

    Furthermore, the computational load of findNearestNode can be reduced through a variety of software optimizations. A simple one used here to demonstrate that fact is storing nodes in seperate ``buckets,'' sorted by their $x$ value. By using only two buckets, the execution time of findNearestNode fell drastically. Figure \ref{subfig:32x32x32rrt2} shows edge collision detection accounting for over 95\% of execution time for $100 \leq K \leq 10000$. This is consistent with the profiling results of \ac{RRT} in prior work\cite{Bialkowski2011}.

    \subsubsection*{Conclusion}
        From the above data, it was identified that, as prior work suggested, edge collision detection shows the greatest promise for potential speedup through specialized hardware. The next chapter details the process of designing and building this hardware.
        \todo[inline]{Add simulations to determine correct K}

    % @Author: AnthonyKenny98
% @Date:   2020-02-29 17:30:44
% @Last Modified by:   AnthonyKenny98
% @Last Modified time: 2020-03-01 08:08:09
\begin{figure}[H]
\begin{center}

    % Subfigure A
    \begin{subfigure}{\textwidth}
    \begin{center}
    \includegraphics[width=\linewidth,height=0.3\paperheight]{chapters/chapter2/img/profiling/16x16x16/performance2.png}
    \caption{16x16x16 Configuration Space}
    \label{subfig:16x16x16rrt2}
    \end{center}
    \end{subfigure}
    % 
    % Subfigure B
    \begin{subfigure}{\textwidth}
    \begin{center}
    \includegraphics[width=\linewidth,height=0.3\paperheight]{chapters/chapter2/img/profiling/32x32x32/performance2.png}
    \caption{32x32x32 Configuration Space}
    \label{subfig:32x32x32rrt2}
    \end{center}
    \end{subfigure}
    % Caption and Label
    \caption{\ac{RRT} Functions Exectution Time, with Bucket Optimization}
    \label{fig:rrt_profiling2}
\end{center}
\end{figure}


% Chapter 3
\chapter{Motion Planning in Hardware}
    \label{chap:MotionPlanningInHardware}
    % @Author: AnthonyKenny98
% @Date:   2020-02-26 08:22:07
% @Last Modified by:   AnthonyKenny98
% @Last Modified time: 2020-04-08 10:40:29

The second objective of this thesis was to design and implement a functional hardware unit that accelerates the execution of \gls{RRT} in \gls{3D}. With the bottleneck function having been identified in Chapter 2 as edge collision detection, Chapter 3 details the specification, design, implementation, and analysis of a hardware unit the implements the edge collision function.

\section{Defining the Collision Detection Unit}
    % @Author: AnthonyKenny98
% @Date:   2020-04-08 09:52:59
% @Last Modified by:   AnthonyKenny98
% @Last Modified time: 2020-04-08 09:53:43
\subsection{Edge Collision Function}
\label{subsection:EdgeCollisionFunction}
    To briefly examine the edge collision detection function in general terms; Given an edge $e$, \gls{RRT} finds where $e$ intersects with grids in the \gls{OGM}. If any of the grids it intersects with are ``occupied'', a collision is returned. This is shown in Figure \ref{fig:edge_collision_process} on Page \pageref{fig:edge_collision_process}. 

    To see why it is so computationally intense to calculate intersections between a segment and grids, it must be understood that it is a fairly involved geometric process. Figure \ref{fig:edge_collision_planes} on Page \pageref{fig:edge_collision_planes} shows how grid intersections are detected by computing the where the segment intersects certain \textbf{\glspl{axis-oriented plane}}.  

    \subsubsection{Time Complexity}
        With the steps of the edge collision algorithm understood, (explained graphically in Figure \ref{fig:edge_collision_planes}, algorithm included in Appendix \ref{section:rrt_appendix_function_impl}, its \Gls{time complexity} may be quantified. 
        For an edge $e$ of maximum length $\epsilon$, it must check for intersections with $\epsilon \times \epsilon \times \epsilon$ grids. (i.e the only grids that are checked are the ones that the $e$ could possibly intersect with). 
        It first iterates through the three dimensions of \glspl{axis-oriented plane} ($xy$, $xz$, and $yz$). This is a constant of 3. 
        Within each of these dimensions, it must iterate through $\epsilon$ planes. This makes its time complexity $O(3\epsilon)$. 
        \todo{Appendix Reference?}% (Appendix C.[TODO] shows an $O(\epsilon^2)$ method upon which this is an improvement). 

        % @Author: AnthonyKenny98
% @Date:   2020-04-07 11:00:50
% @Last Modified by:   AnthonyKenny98
% @Last Modified time: 2020-04-07 12:40:47
\begin{figure}[H]
\begin{centering}
\begin{tabular}{cc}
    \begin{subfigure}{0.47\linewidth}
    \includegraphics[width=\linewidth]{chapters/chapter3/img/edge_collision_planes_a.png}
    \caption{}
    \label{fig:edge_collision_planes_a}
    \end{subfigure} &

    \begin{subfigure}{0.47\linewidth}
    \includegraphics[width=\linewidth]{chapters/chapter3/img/edge_collision_planes_b.png}
    \caption{}
    \label{fig:edge_collision_planes_b}
    \end{subfigure} \\

    \begin{subfigure}{0.47\linewidth}
    \includegraphics[width=\linewidth]{chapters/chapter3/img/edge_collision_planes_c.png}
    \caption{}
    \label{fig:edge_collision_planes_c}
    \end{subfigure} &

    \begin{subfigure}{0.47\linewidth}
    \includegraphics[width=\linewidth]{chapters/chapter3/img/edge_collision_planes_d.png}
    \caption{}
    \label{fig:edge_collision_planes_d}
    \end{subfigure} \\
\end{tabular}

\mycaption{Detecting Grid Intersections by Finding Intersections with Axis Oriented Planes}{\\ Consider an edge $e$ that spans the width of a $2\times 2\times 2$ grid map, as shown in Figure \ref{fig:edge_collision_planes} (do not consider if the grids are occupied, this is just to determine which of them the edge intersects). Just by eyeballing Figure \ref{fig:edge_collision_planes_a} it seems odd that so many of the grids have been intersected (denoted by red shading) by the yellow edge. The algorithm executed by checking one set of \glspl{axis-oriented plane} at a time. Figure \ref{fig:edge_collision_planes_a} shows how the $xy$-oriented planes are checked for 3 different values of $z$ (going into the page), finding two intersection points. There is only one intersection for the $xz$ oriented planes (\ref{fig:edge_collision_planes_c}). In Figure \ref{fig:edge_collision_planes_d}, the segment intersects 2 grids in the second $yz$-oriented plane, and one in the third plane. The intersected grids are thus any grid where a point-of-intersection falls on its face.}
\label{fig:edge_collision_planes}
\end{centering}
\end{figure}

        % @Author: AnthonyKenny98
% @Date:   2020-04-06 17:52:43
% @Last Modified by:   AnthonyKenny98
% @Last Modified time: 2020-04-07 12:37:57

\begin{figure}[H]
\begin{centering}
\includegraphics[width=\linewidth]{chapters/chapter3/img/edge_collision_process.png}
\mycaption{Edge Collision Computation Process}{. Most of the computational load is found in determining the grids with which an edge intersects (as seen on the left of the figure). Once these grids have been found, it is very simple (and fast) to lookup if these grids are occupied in the \gls{OGM}}
\label{fig:edge_collision_process}
\end{centering}
\end{figure}
   
\subsection{Technical Specifications}
\label{subsection:HoneyBeeSpecs}

    \subsubsection{Performance Specifications}
        When accelerating motion planning algorithms, it is often difficult to quantify a goal for how much faster one would like the function to run - the answer is usually ``as fast as possible!'' For this thesis, the performance specification was set to be the edge collision function to run fast enough such that it was no longer the bottleneck function. This translated to a desired speedup of about 3 times (when compared to benchmark performance of a typical CPU). Table \ref{table:edg_col_performance_specs} quantifies this in terms of latency and throughput.

        % @Author: AnthonyKenny98
% @Date:   2020-04-07 13:28:39
% @Last Modified by:   AnthonyKenny98
% @Last Modified time: 2020-04-08 16:26:52
\begin{table}[H]
\begin{centering}
\begin{tabular}{|c|c|c|}
\hline
\textbf{Metric}     &   \textbf{Benchmark CPU$^*$}   & \textbf{Accelerated} \\
\hline
Latency ($\mu$seconds/edge)  &   2.6  & 0.9 \\
\hline
Throughput (edges/second)  &   384,615  & 1,111,111 \\
\hline
\end{tabular}
\mycaption{Performance Specifications for Edge Collision Detection Unit}{. $^*$Benchmark CPU is an Intel 3.1 GHz i7 Dual Core processor, typical of a laptop computer.}
\label{table:edg_col_performance_specs}
\end{centering}
\end{table}

\newpage

\subsubsection{Area Specifications}
    Generally, an inverse relationship exists between latency and area. While it may be possible to make the unit much faster than the latency specification, this may become prohibitive with regards to the amount of area on chip it would occupy. It was decided to limit the area to that which would fit on an \gls{FPGA} typical in drone applications (those of the Kintex-7 Low Voltage family were chosen, but there are many possible options.) 

    Logic area on an FPGA is largely determined by \glspl{LUT}. \glsdesc{LUT_g} As such, the upper bound on area was set at 274,080 \glspl{LUT}.

    \subsubsection{Interface Specifications}
    \label{subsection:HoneyBeeTechSpechs}
        As shown in Figure \ref{fig:edge_collision_process}, the computationally intensive part of the process of edge collision detection is finding points of intersection between an edge and the grids of the map. Comparing this result to an \gls{OGM} is simple and fast. Therefore, it was decided that the hardware unit would simply take an edge and determine the grids with which it intersects. Whether the edge intersects a given grid can be represented as a binary $\{0,1\}$, and thus a the intersections found in a $\epsilon \times \epsilon\times\epsilon$ gridspace can be represented as an $\epsilon^3$ sequence of binary values.
        Table \ref{table:edg_col_interface_specs} outlines the required interface specifications for the functional unit.
        % @Author: AnthonyKenny98
% @Date:   2020-03-01 14:11:33
% @Last Modified by:   AnthonyKenny98
% @Last Modified time: 2020-04-07 14:09:51
\begin{table}[H]
\begin{center}
\begin{tabular}{|p{.2\linewidth}|p{.74\linewidth}|}
    \hline
    \textbf{Element}             & \textbf{Description/Justification} \\
    \hline
    \multicolumn{2}{|c|}{Constraints} \\
    \hline
    Length $\epsilon$  & $\epsilon$ defines the max edge length. The space being checked and the output sequence has the dimensions $\epsilon\times\epsilon\times\epsilon$ \\
    \hline
    \multicolumn{2}{|c|}{Inputs} \\
    \hline
    Edge $e$  & An Edge $e$ defined for a \gls{3D} \gls{configuration} space by two points $\{p1, p2\}$, each defined by a set of \gls{3D} coordinates $\{x,y,z\}$.\\
    \hline
    Control Inputs & The functional unit must have ports for control signals: clock, reset, start. These are required for adding the unit to a processor. \\
    \hline
    \multicolumn{2}{|c|}{Outputs} \\ 
    \hline
    Return Value & $\epsilon^3$ bit sequence: 1 if collides with grid at that index, 0 otherwise.\\
    \hline
    Control Outputs & Output ports for control signals: idle, done, ready. These are required for adding the unit to a processor. \\
    \hline
\end{tabular}
\mycaption{Interface Specifications for Edge Collision Detection Unit}{}
\label{table:edg_col_interface_specs}
\end{center}
\end{table}

\newpage

\section{HoneyBee}
    % @Author: AnthonyKenny98
% @Date:   2020-04-08 09:56:57
% @Last Modified by:   AnthonyKenny98
% @Last Modified time: 2020-04-10 15:24:37

% @Author: AnthonyKenny98
% @Date:   2020-04-06 17:19:39
% @Last Modified by:   AnthonyKenny98
% @Last Modified time: 2020-04-06 17:29:38
\begin{figure}[H]
\begin{centering}
\includegraphics[width=0.8\linewidth]{chapters/chapter3/img/MegalongHoneyBee.png}
\mycaption{Megalong Park Honey Bee Pollinating a Weeping Cherry Blossom}{. Photographed by Emma Kenny in the Southern Highlands of New South Wales, Australia}
\end{centering}
\end{figure}

The honey bee, \textit{Apis mellifera}, has long been renowned for its tireless work ethic. However, the it is rarely given credit for its remarkable navigation and collision avoidance strategies during flight. Recent research\cite{Menzel2005} suggests that honey bees, interestingly enough, explore their workspace randomly in order to find paths from their hive to sources of pollen. Sound familiar? As such, it is quite appropriate that this functional unit, designed to work tirelessly, rapidly and efficiently to execute collision detection computations for robot motion planning, was named \textbf{HoneyBee}. \\
\bigskip

HoneyBee is a hardware unit that will eventually be incorporated into a processor, demonstrated in Figure \ref{fig:honeybee_in_processor}. In Chapter 4, the HoneyBee unit is implemented in a simple RISC-V processor and invoked using custom RISC-V instructions. For now, however, consider HoneyBee as a standalone unit that computes the grids with which an edge collides. Its resulting output can be compared to an \gls{OGM}, as explained in section \ref{subsection:EdgeCollisionFunction}.

% @Author: AnthonyKenny98
% @Date:   2020-04-08 11:15:41
% @Last Modified by:   AnthonyKenny98
% @Last Modified time: 2020-04-08 11:20:08
\begin{figure}[H]
\begin{centering}
% \includegraphics[width=\linewidth]{name}
\missingfigure[figwidth=\linewidth]{HoneyBee in a Processor}
\caption{}
\label{fig:honeybee_in_processor}
\end{centering}
\end{figure}

\subsection{HoneyBee Interface Design}
    The interface for the HoneyBee functional unit, following on from the interface specifications outlined in Section \ref{subsection:HoneyBeeTechSpechs} can be simply represented by Figure \ref{fig:honeybee_interface_simple}.

    % @Author: AnthonyKenny98
% @Date:   2020-04-08 10:23:56
% @Last Modified by:   AnthonyKenny98
% @Last Modified time: 2020-04-08 10:44:54
\begin{figure}[H]
\begin{centering}
\includegraphics[width=0.6\linewidth]{chapters/chapter3/img/honeybee_interface_simple.png}
\mycaption{General Overview of HoneyBee Interface}{. The functional unit takes an edge $e$, defined by two points $p_1$ and $p_2$, as an input, and outputs a series of collisions. These collisions describe which grids an edge intersects. Its control interface allows for communication with a processors main control unit.}
\label{fig:honeybee_interface_simple}
\end{centering}
\end{figure}
    
    However, when designing hardware (the method of doing so is described in section \ref{section:honeybee_implementation}), how these inputs and outputs are implemented must be considered at the \gls{bit} level. Figure \ref{fig:honeybee_interface_detail} shows all input, output, and control ports, and their \glspl{bit-width}. 
    
    % @Author: AnthonyKenny98
% @Date:   2020-04-08 10:33:04
% @Last Modified by:   AnthonyKenny98
% @Last Modified time: 2020-04-10 15:09:42
\begin{figure}[H]
\begin{centering}
\includegraphics[width=0.8\linewidth]{chapters/chapter3/img/honeybee_interface_detail.png}
\mycaption{Port Diagram of HoneyBee Interface}{. The edge input is represented by 6 32-bit floats, following the IEEE 754 Single Precision 32-bit protocol, with each float representing one of its coordinate points. The output sequence of collisions is an $e^3$-bit sequence, with each bit in the sequence representing one of the grids that was checked for intersections. It has input control signals for start and reset, and output control signals for done, idle, and ready. These control signals make up the necessary signals for a handshake protocol between HoneyBee and a processor's main controller.}
\label{fig:honeybee_interface_detail}
\end{centering}
\end{figure}

    \subsubsection{Inputs}
        The inputs to HoneyBee collectively describe a single edge. This is done with 6 32-bit floating point numbers. How floating points (non-integer numbers) are represented in binary is defined by the \gls{IEEE754}. How this is actually represented is not neccesary to understand, but is explained in Appendix \ref{section:honeybee_appendix_ieee}. The important point is that the input edge is determined by 6 32-bit coordinate points.

    \subsubsection{Output}
        HoneyBee outputs a sequence of ``collision-bits,'' with each bit in the sequence representing any collisions between the input edge with its corresponding grid. How this sequence of bits is mapped to a \gls{3D} grid-map is explained in Appendix \ref{section:honeybee_appendix_mapping}. It is important to note that in the design and implementation of HoneyBee, the length of this sequence was parameterized to be variable, corresponding to a variable value of $\epsilon$. Recall that the optimal edge collision algorithm only checked $\epsilon^3$ grids. HoneyBee, as well, only checks the grids with which the edge could possibly intersect.

        Since the number of grids being checked is parameterized, so must the number of collision-bits. This is demonstrated in Figure \ref{fig:honeybee_epsilon_grids}.
        % @Author: AnthonyKenny98
% @Date:   2020-04-08 11:48:44
% @Last Modified by:   AnthonyKenny98
% @Last Modified time: 2020-04-10 15:13:02
\begin{figure}
\begin{centering}
\begin{tabular}{c c}

    \begin{subfigure}{0.4\linewidth}
    \includegraphics[width=\linewidth]{chapters/chapter3/img/honeybee_epsilon_grids_a.png}
    \caption{For $\epsilon=2$, an output collision-bit sequence of length 8 is required.}
    \end{subfigure} & 

    \begin{subfigure}{0.4\linewidth}
    \includegraphics[width=\linewidth]{chapters/chapter3/img/honeybee_epsilon_grids_b.png}
    \caption{For $\epsilon=4$, an output collision-bit sequence of length 64 is required.}
    \end{subfigure} \\    

\end{tabular}
\mycaption{The Impact of $\epsilon$ on the Length of the Bit-Collision Sequence}{}
\label{fig:honeybee_epsilon_grids}
\end{centering}
\end{figure}

        \textbf{Note:} The output \gls{bit-width} is \textit{parameterized} not \textit{variable}. Upon synthesis (building) of HoneyBee, the output \gls{bit-width} is set at a constant value. Different syntheses may have different output \glspl{bit-width}. When the time comes to add HoneyBee to a processor, it is synthesized with a certain \gls{bit-width}.

    \subsubsection{Control Interface}
        The control interface is designed to give HoneyBee the ability to be included in a processor, and implements a commonly used ``handshake'' protocol between HoneyBee and the control unit of the processor in which it resides. Put simply, this is a method that allows the control unit to tell HoneyBee when to start executing the computation, and for HoneyBee to tell the control unit when it has finished its computation and the output value is ready. This is explained in detail in Appendix \ref{section:honeybee_appendix_handshake}. The control interface also has a clock and reset port.


\subsection{HoneyBee Implementation}
\label{section:honeybee_implementation}
    \subsubsection{Hardware Description Languages}
        Designing computers and their constituent parts has come a long way from its arduous beginnings. ``Victory'', the enigma-breaking machine designed by Alan Turing at Bletchley Park during World War II, was a large electro-mechanical computer made up of storage wheels, electromagnetic relays, and rotary switches, assembled by hand.\cite{ChoiceReviews2006} So too was ``Mark I'', the 816 cubic feet computer designed by Harvard University's Dr. Howard Aiken, which, on March 1944, computed the viability of implosion for detonating the atomic bomb.\cite{Elsabbagh2019}
        
        % @Author: AnthonyKenny98
% @Date:   2020-04-08 13:30:24
% @Last Modified by:   AnthonyKenny98
% @Last Modified time: 2020-04-08 13:54:06
\begin{figure}
\begin{centering}
\includegraphics[width=0.7\linewidth]{chapters/chapter3/img/mark1.jpeg}
\mycaption{The Harvard Mark I Computer}{, photographed in the Harvard Science Center in April 2014 (Source: Harvard University)}
\end{centering}
\end{figure}

        Computers nowadays measure in the order of millimeters rather than meters. What's more, they are now ``built'' in software, using a \glsfirst{HDL}. \glspl{HDL} are a family of computer programming languages that are used to specify the function of electronic circuits. Tools allow for simulation of such circuits to verify design correctness and performance. Modules defined in \glspl{HDL} may then be synthesized for a type of integrated circuit called a \glsfirst{FPGA}. This \gls{FPGA}, ``programmed'' in \gls{HDL} code to behave in a certain way, can then serve the purpose of a processor or other functional processing unit. Figure \ref{fig:hdl_to_fpga} demonstrates this process.

        % @Author: AnthonyKenny98
% @Date:   2020-04-08 13:37:54
% @Last Modified by:   AnthonyKenny98
% @Last Modified time: 2020-04-10 15:14:42
\begin{figure}[H]
\begin{centering}
\includegraphics[width=\linewidth]{chapters/chapter3/img/hdl_to_fpga.png}
\mycaption{The Hardware Development Process}{, Defining Hardware Units in Hardware Description Languages for FPGAs.}
\label{fig:hdl_to_fpga}
\end{centering}
\end{figure}

        HoneyBee was implemented eventually in an \gls{HDL} called Verilog. However, no Verilog for HoneyBee was ever explicitly written by a human. It was generated by a tool called High-Level Synthesis.

    \subsubsection{High Level Synthesis}
        \gls{HLS} is an automated hardware design process that takes design files (written in high-level languages, such as C, C++ or SystemC) specifying the algorithmic function of a piece of hardware, interprets those files, and creates digital hardware designs that execute this function. In short, it effectively translates programming languages into hardware description languages. Some key advantages of using HLS are speed and verification. It is much faster and easier to define functionality in C than it is in a \gls{HDL} such as Verilog, and thus design iterations are faster. It is also much simpler to verify one's design, as the functional units can be put through test benches written in C. \\

        The most important benefit of using \gls{HLS}, however, is the ability to use ``pragmas.'' These are directives given to the \gls{HLS} tool that tell it what optimizations to use when translating C code into an \gls{HDL}. This allows the same funtionality to be synthesized in many different ways, optimizing the synthesis for speed, area, memory, etc. As such, this hardware development process allows developers to experiment quickly with different ways to implement the same functionality. This is demonstrated graphically by Figure \ref{fig:hls_to_fpga} on Page \pageref{fig:hls_to_fpga}.

        % @Author: AnthonyKenny98
% @Date:   2020-04-08 15:17:43
% @Last Modified by:   AnthonyKenny98
% @Last Modified time: 2020-04-08 16:02:17
\begin{figure}
\begin{centering}
\begin{tabular}{c}

\begin{subfigure}{0.97\linewidth}
\includegraphics[width=\linewidth]{chapters/chapter3/img/hls_to_fpga_a.png}
\caption{Hardware Optimization without HLS}
\label{fig:hls_to_fpga_a}
\end{subfigure} \\

\begin{subfigure}{0.97\linewidth}
\includegraphics[width=\linewidth]{chapters/chapter3/img/hls_to_fpga_b.png}
\caption{Hardware Optimization with HLS}
\label{fig:hls_to_fpga_b}
\end{subfigure}

\end{tabular}
\mycaption{Hardware Optimization Process}{. Figure \ref{fig:hls_to_fpga_a} shows how, without using HLS, the hardware developer must write multiple different implementation of the same functionality to acheive different performance. Figure \ref{fig:hls_to_fpga_b} shows how, when using HLS, the hardware developer only must write one implementation (in a higher-level language like C), and then use ``pragmas'' to create different implementations (with different optimizations) for synthesis.}
\label{fig:hls_to_fpga}
\end{centering}
\end{figure}

    \subsubsection{HoneyBee-A Synthesis}
        With the functionality of HoneyBee implemented in C, the first optimization iteration (designated \gls{HB-A}) was synthesized. \gls{HB-A} had no pragmas, and was merely a basic hardware implementation of the defined functionality. \gls{HB-A} (and all subsequent iterations) synthesized correctly and satisfied the interface specifications (See Appendix \ref{section:honeybee_appendix_synthesis_report} for technical interface report of synthesis). Table \ref{table:HBA_results} shows the result of \gls{HB-A}'s synthesis compared with its performance and area specifications. Obviously, the synthesis is well below the area limit, but nowhere near the specifed performance metrics. This is where the beauty of High-Level Synthesis optimization comes in.

        % @Author: AnthonyKenny98
% @Date:   2020-04-08 15:34:16
% @Last Modified by:   AnthonyKenny98
% @Last Modified time: 2020-04-08 15:49:09
\begin{table}[H]
\begin{center}
\begin{tabular}{|c|c|c|}
    \hline
    \textbf{Metric}             & \textbf{Specification} & \textbf{Synthesis Result} \\
    \hline
    Latency ($\mu$seconds/edge)  &   0.9  & \textcolor{myred}{6.66} \\
    \hline
    Throughput (edges/second)  &   1,111,111 & \textcolor{myred}{150,150} \\
    \hline
    FPGA Area \glspl{LUT} & 274,080  & \textcolor{mygreen}{10,593} \\
    \hline
\end{tabular}
\mycaption{Synthesis Results for HB-A with $\epsilon = 4$}{}
\label{table:HBA_results}
\end{center}
\end{table}

\newpage
\subsection{HoneyBee Acceleration}
    This section steps through the process of using \gls{HLS} \glspl{pragma} in 2 major optimization iterations, \gls{HB-B} and \gls{HB-C} and how these iterations compared to benchmark and specified performance.

    \subsubsection{Benchmarking}
        The benchmark performance was based on a Dual-Core Intel 3.1GHz i7 processor. In an ideal world, this processor would have been chosen after a rigourous process of determining the most suitable benchmark. In reality, this processor was chosen because it was the one found in the computer running the simulations (an early 2015 MacBook Pro). Nevertheless, it serves as a suitable benchmark, demonstrating performance typical of general purpose CPUs.

        Examining latency, the benchmark was set at the average execution time for the benchmark CPU to compute 1 edge. Tests were run and averaged over 1000 trials. The average latency was 2.6 $\mu$seconds, with a standard deviation of 0.1 $\mu$seconds. This is shown, along with the performance of \gls{HB-A} in Figure \ref{fig:benchmark_1}.

        % @Author: AnthonyKenny98
% @Date:   2020-04-08 16:42:38
% @Last Modified by:   AnthonyKenny98
% @Last Modified time: 2020-04-08 17:12:56
\begin{figure}[H]
\begin{centering}
\includegraphics[width=\linewidth]{chapters/chapter3/img/benchmark_1.png}
\mycaption{HB-A Performance Against Benchmark CPU}{. Standard deviation of CPU average performance is shown by width of the line for easier interpretation.}
\label{fig:benchmark_1}
\end{centering}
\end{figure}

    \subsubsection{HoneyBee-B}
        The first step in accelerating HoneyBee was taking advantage of the inherent parrallelism available to the algorithm. Recall that the edge collision algorithm checks the $xy$-oriented planes, the $xz$-oriented planes, and finally the $yz$-oriented planes. These operations are completely independent, and can thus be performed simultaneously. Figure \ref{fig:hbb_timing} shows how the process of checking each set of planes was done sequentially in \gls{HB-A}, but are executed in parallel in \gls{HB-B}.

        % @Author: AnthonyKenny98
% @Date:   2020-04-08 18:12:53
% @Last Modified by:   AnthonyKenny98
% @Last Modified time: 2020-04-08 18:42:07
\begin{figure}[H]
\begin{centering}
\begin{tabular}{c}

\begin{subfigure}{\linewidth}
\includegraphics[width=\linewidth]{chapters/chapter3/img/timing1.png}
\caption{HoneyBee-A Timing Diagram. \texttt{Check\_Planes} executed sequentially.}
\label{fig:hbb_timing_a}
\end{subfigure} \\

\begin{subfigure}{\linewidth}
\includegraphics[width=\linewidth]{chapters/chapter3/img/timing2.png}
\caption{HoneyBee-B Timing Diagram. \texttt{Check\_Planes} executed in parallel. Different colors are to represent slightly different implementations of the same function.}
\label{fig:hbb_timing_b}
\end{subfigure} \\

\end{tabular}
\mycaption{Timing Diagrams Showing Parallelization in HoneyBee-B}{. Note these are generalizations for easy explanation of the concept of hardware parallelization. Full timing reports for each HoneyBee iteration can be found in Appendix \ref{section:honeybee_appendix_timing_reports}}
\label{fig:hbb_timing}
\end{centering}
\end{figure}

        The timing diagram does not only show the use of parallelism, it also shows how this is actually implemented in hardware. In \gls{HB-A} (Figure \ref{fig:hbb_timing_a}), there is a single module named \texttt{Check\_Plane}. Since there is only one of them, it must calculate intersections with each set of planes sequentially. On the other hand, in \gls{HB-B} (Figure \ref{fig:hbb_timing_b}), there are three separate instances of this module (\texttt{Check\_Planes\_XY}, \texttt{Check\_Planes\_XZ}, and \texttt{Check\_Planes\_YZ}), allowing HoneyBee to execute computation on all three sets of planes in parallel. 

        Moreover, notice that in \gls{HB-B}, execution time of each of these three instances is shorter than that of the single instance in \gls{HB-A}. When a single module instance is used for different purposes (in this case, checking the $xy$, $xz$, and $yz$ oriented planes), it has some variability. To control this variability, it must execute a certain amount of control logic at the beginning of the function. On the other hand, when there are seperate instances of the module, what was once variable can now be made constant. As a result, each instance can be slightly specialized and the control logic eliminated. In this case, a general \texttt{Check\_Planes} module was replaced with the three specialized \texttt{Check\_Planes\_XY}, \texttt{Check\_Planes\_XZ}, and \texttt{Check\_Planes\_YZ}. Each of these has less variability, thus less control logic, and therefore a slightly faster overall latency. Figure \ref{fig:hbb_timing} shows how theoretically, this should result in a reduction in overall latency of more than 3 times.\\

        Comparing HoneyBee-B against our benchmark, the success of this optimization can be seen. Figure \ref{fig:benchmark_2} shows the performance of multiple variants of \gls{HB-B} in yellow. These variants were the result of experimenting with slightly different \glspl{pragma}, but all fell in roughly the same area. Appendix \ref{section:honeybee_appendix_hbb_variants} lists the details of each \gls{HB-B} variant. \\
        HB-B3 showed the best performance/area relationship. It was of acceptable area and had latency marginally lower than the benchmark, but was not as fast as the defined specifications. Table \ref{table:HBB_results} shows the result of HB-B3's synthesis compared with its performance and area specifications.

        % @Author: AnthonyKenny98
% @Date:   2020-04-08 15:34:16
% @Last Modified by:   AnthonyKenny98
% @Last Modified time: 2020-04-08 19:38:23
\begin{table}[H]
\begin{center}
\begin{tabular}{|c|c|c|}
    \hline
    \textbf{Metric}             & \textbf{Specification} & \textbf{Synthesis Result} \\
    \hline
    Latency ($\mu$seconds/edge)  &   0.9  & \textcolor{myred}{2.05} \\
    \hline
    Throughput (edges/second)  &   1,111,111 & \textcolor{myred}{487,805} \\
    \hline
    FPGA Area \glspl{LUT} & 274,080  & \textcolor{mygreen}{26,524} \\
    \hline
\end{tabular}
\mycaption{Synthesis Results for HB-B3 with $\epsilon = 4$}{}
\label{table:HBB_results}
\end{center}
\end{table}

        % @Author: AnthonyKenny98
% @Date:   2020-04-08 17:54:49
% @Last Modified by:   AnthonyKenny98
% @Last Modified time: 2020-04-08 19:08:39
\begin{figure}[H]
\begin{centering}
\includegraphics[width=\linewidth]{chapters/chapter3/img/benchmark_2.png}
\mycaption{HB-B Performance Against Benchmark CPU}{\\ Variants of HB-B shown in Yellow.}
\label{fig:benchmark_2}
\end{centering}
\end{figure}

    \subsubsection{HoneyBee-C}
        A similar concept was applied to optimize the computation of each set of planes. Consider synthesis of HoneyBee for $\epsilon = 4$. HoneyBee will need to compute intersections with 4 $xy$-oriented planes, 4 $xz$-oriented planes, and 4 $yz$-oriented planes. \gls{HB-B} computes each set of 4 planes simultaneously, but the computing intersections with each of the 4 planes in one orientation are also independent operations. As such, they can also be parallelized with the instantiation of more hardware modules. Figure \ref{fig:hbc_timing} shows the timing diagrams for the \texttt{Check\_Planes\_XY} module that was shown in the last set of timing diagrams.

        % @Author: AnthonyKenny98
% @Date:   2020-04-08 18:12:53
% @Last Modified by:   AnthonyKenny98
% @Last Modified time: 2020-04-10 16:29:06
\begin{figure}[H]
\begin{centering}
\begin{tabular}{c}

\begin{subfigure}{0.97\textwidth}
\includegraphics[width=\linewidth]{chapters/chapter3/img/timing3.png}
\caption{HoneyBee-B Timing Diagram for \texttt{Check\_Planes\_XY}. One instance of \texttt{Check\_Planes\_XY} module executed sequentially 4 times ($\epsilon=4$.}
\label{fig:hbc_timing_a}
\end{subfigure} \\

\begin{subfigure}{0.97\textwidth}
\includegraphics[width=\linewidth]{chapters/chapter3/img/timing4.png}
\caption{HoneyBee-C Timing Diagram. 4 instances of \texttt{Check\_Planes\_XY} module executing in parallel.}
\label{fig:hbc_timing_b}
\end{subfigure} \\

\end{tabular}
\mycaption{Timing Diagrams Showing Parallelization in HoneyBee-C}{. Again, these are simplified versions of timing analysis for easy explanation of the concept of hardware parallelization.}
\label{fig:hbc_timing}
\end{centering}
\end{figure}

        Just focussing on the computation of the $xy$-oriented planes, Figure \ref{fig:hbc_timing} shows how \gls{HB-B} has only one instance of the module execute 4 times sequentially. \gls{HB-C}, on the other hand, implements greater parallelism by instantiating 4 module instances to execute sequentially. \gls{HB-C}, as a result, exceeded the performance benchmark and was of an acceptable area, as shown in Table \ref{table:HBC_results} and Figure \ref{fig:benchmark_3}.

        % @Author: AnthonyKenny98
% @Date:   2020-04-08 15:34:16
% @Last Modified by:   AnthonyKenny98
% @Last Modified time: 2020-04-08 19:39:39
\begin{table}[H]
\begin{center}
\begin{tabular}{|c|c|c|}
    \hline
    \textbf{Metric}             & \textbf{Specification} & \textbf{Synthesis Result} \\
    \hline
    Latency ($\mu$seconds/edge)  &   0.9  & \textcolor{mygreen}{0.53} \\
    \hline
    Throughput (edges/second)  &   1,111,111 & \textcolor{mygreen}{1,886,792} \\
    \hline
    FPGA Area \glspl{LUT} & 274,080  & \textcolor{mygreen}{185,663} \\
    \hline
\end{tabular}
\mycaption{Synthesis Results for HB-C with $\epsilon = 4$}{}
\label{table:HBC_results}
\end{center}
\end{table}

        % @Author: AnthonyKenny98
% @Date:   2020-04-08 17:54:49
% @Last Modified by:   AnthonyKenny98
% @Last Modified time: 2020-04-10 15:26:35
\begin{figure}[H]
\begin{centering}
\includegraphics[width=\linewidth]{chapters/chapter3/img/benchmark_3.png}
\mycaption{HoneyBee-C Performance Against Benchmark CPU}{\\ HB-B variants shown in yellow, HB-C in green.}
\label{fig:benchmark_3}
\end{centering}
\end{figure}

% Chapter 4
\chapter{RISC-V Processor}
    \label{chap:RiscvProcessor}
    % @Author: AnthonyKenny98
% @Date:   2020-02-26 09:00:23
% @Last Modified by:   AnthonyKenny98
% @Last Modified time: 2020-03-01 19:53:23

\section{Introduction to the Reduced Instruction Set Computer}
    % @Author: AnthonyKenny98
% @Date:   2020-03-01 19:39:42
% @Last Modified by:   AnthonyKenny98
% @Last Modified time: 2020-04-03 13:33:06
    
\subsection{Instruction Set Architecture}
    An Instruction Set Architecture can be thought of as an abstract model of a computer. On a broad level, it defines the data types, memory model, and registers of a computer, along with the instructions that it can execute. \\
    It can also be thought as a ``contract'' between hardware and software developers. It is the promise made that the hardware will be able to execute all instructions defined in the \gls{ISA}, and the limitation that software must be compiled into that set of instructions.

% \begin{figure}[H]
% \begin{center}
%     \begin{subfigure}{0.4\textwidth}
%         \begin{lstlisting}[style=CStyle]
%     #include <stdio.h>
%     int main(void) {
%         c = a + b; 
%     } 
%         \end{lstlisting}
%         \caption{C Function}
%     \end{subfigure} 
%     \ \ \ \ \
%     \begin{subfigure}{0.4\textwidth}
%         \begin{lstlisting} %[style=ASM]
%     add $x3, $x1, $x2 
%         \end{lstlisting}
%         \caption{Assembly Implementation}
%     \end{subfigure}

% \caption{Example of C function and its equivalent assembly code}
% \end{center}
% \end{figure}

\subsection{RISC Processor Design}
    
    % @Author: AnthonyKenny98
% @Date:   2020-03-01 20:09:46
% @Last Modified by:   AnthonyKenny98
% @Last Modified time: 2020-04-09 11:51:30
\begin{figure}[H]
    \includegraphics[width=\linewidth]{chapters/chapter4/img/RISC-Datapath.png}
    \mycaption{5-Stage \gls{RISC} Datapath}{. Combinational logic stages are shown in yellow, and registers between stages are shown in green.}
    \label{fig:RISC-Datapath}
\end{figure}

\section{RISC-V ISA}
    % @Author: AnthonyKenny98
% @Date:   2020-03-01 18:51:24
% @Last Modified by:   AnthonyKenny98
% @Last Modified time: 2020-03-01 19:39:08
\subsection{RV32I}

        The following is an excerpt from the RISC-V Specification, outlining the RV32I base integer instruction set \cite{Waterman2019}
        \begin{quote}{}
            \small{RV32I was designed to be sufficient to form a compiler target and to support modern operating system environments. The ISA was also designed to reduce the hardware required in a minimal implementation. RV32I contains 40 unique instructions, though a simple implementation might \dots [reduce] base instruction count to 38 total. RV32I can emulate almost any other ISA extension \dots \\
            Subsets of the base integer ISA might be useful for pedagogical purposes, but the base has been defined such that there should be little incentive to subset a real hardware implementation \dots}
        \end{quote}

        \subsubsection*{Registers}
        RV32I defines 32 unprivileged registers, each 32 bits wide. They are designated \texttt{x0-x31}, where \texttt{x0} is a hard-wired value of $0$, and registers \texttt{x1-x31} hole values that various instructions use. RISC-V uses the load-store method, meaning that all operations perform on two registers or a register and an immediate, rather than performing operations directly on memory addresses. In addition, the 33rd unprivileged register is the program counter \texttt{pc}. Table \ref{table:rv32i_reg} shows the register state for the RV32I Base Integer Instruction Set.
        % @Author: AnthonyKenny98
% @Date:   2020-03-01 18:36:35
% @Last Modified by:   AnthonyKenny98
% @Last Modified time: 2020-03-01 19:33:20
\begin{table}[H]
\begin{center}
\begin{tabular}{|p{.15\linewidth}|p{.15\linewidth}|p{.4\linewidth}|}
    \hline
    \textbf{Register}   & \textbf{ABI Name}  & \textbf{Description} \\
    \hline
    \texttt{x0}  & \texttt{zero} & Hard-wired zero \\
    \texttt{x1}& \texttt{ra}& Return address\\
    \texttt{x2}& \texttt{sp} & Stack pointer\\
    \texttt{x3}& \texttt{gp}&Global pointer\\
    \texttt{x4}& \texttt{tp}& Thread pointer\\
    \texttt{x5-7}& \texttt{t0-2}&Temporaries\\
    \texttt{x8}& \texttt{s0/fp}&Saved register/Frame pointer\\
    \texttt{x9}&\texttt{s1} &Saved register\\
    \texttt{x10-11}&\texttt{a0-1}&Function arguments/return values\\
    \texttt{x12-17}&\texttt{a2-7}&Function arguments\\
    \texttt{x18-27}&\texttt{s2-11}&Saved registers\\
    \texttt{x28-31}&\texttt{t3-6}&Temporaries\\
    \hline
    \texttt{pc} & \texttt{pc} & Program counter \\
    \hline
\end{tabular}
\caption{Register State for RV32I Base Instruction Set}
\label{table:rv32i_reg}
\end{center}
\end{table}

        \subsubsection{Instruction Formats}
        Table \ref{table:rv32i_instr_format} demonstrates the format of each different instruction type. 
        % @Author: AnthonyKenny98
% @Date:   2020-03-01 19:13:03
% @Last Modified by:   AnthonyKenny98
% @Last Modified time: 2020-03-01 19:32:00
\begin{table}[H]
\begin{center}
    \includegraphics[draft=false,width=\linewidth]{chapters/chapter4/img/rv32i_instr_format.png}
    \caption{RV32I Base Instruction Formats}
    \label{table:rv32i_instr_format}
\end{center}
\end{table}



    \subsection{Motion Planning Extension}
        \todo[inline, caption={Motion Planning Extension}]{Full description of design of Non standard extension for motion planning. Should follow define, design, build, measure, analyse etc format.}

\section{PhilosophyV}
    \textit{Philosophy IV}, written in 1903 by Mr. Owen Wister of the Class of 1882, recounts the antics of two Harvard students and their last minute attempts to study (or avoid studying) for an exam for which they are hopelessly unprepared. Similarly, this section details the process of my attempt to build a RISC-V processor, by far the most complex part of this Thesis, and a task for which I am unsure of my preparedness. As such, this processor is named Philosphy V; both in reference to the RISC-V ISA for which it is designed, and to the fact that my current situation seems much like a sequel to Mr. Wister's novel.

    \subsection{Baseline Implementation}
        \todo[inline]{Description of Baseline Philosophy V core}
        Figure \ref{fig:philv-core} provides a schematic of the PhilosophyV processor.
        \todo[inline]{Display bigger version of processor.}
        % \newpage
        % \thispagestyle{empty}
        % @Author: AnthonyKenny98
% @Date:   2020-03-01 15:57:31
% @Last Modified by:   AnthonyKenny98
% @Last Modified time: 2020-04-10 06:41:53
\begin{figure}[H]
\begin{centering}
\includegraphics[width=0.7\linewidth]{chapters/chapter4/img/philv-core.png}
\mycaption{RV32I PhilosophyV Schematic}{}
\label{fig:philv-core}
\end{centering}
\end{figure}
 

    \subsection{Implementing HoneyBee}
        \todo[inline]{Process of implementing honeybee into PhilV.}

\section{Performance Analysis}
    \todo[inline]{Comparative Performance Analysis of baseline and extended PhilV Core}

% Conclusion
\chapter{Conclusion}
    \section{Discussion of Results}
        \todo[inline]{Discussion of Results}
    \section{Evaluation of Success}
        \todo[inline]{Evaluation of Success}
    \section{Future Work}
        \todo[inline]{Future Work}

\bibliography{\bibliographyName}
\bibliographystyle{ieeetr}

\begin{appendices}

\chapter{Project Repository}
    \label{appendix:repository}
    % @Author: AnthonyKenny98
% @Date:   2020-02-23 14:22:13
% @Last Modified by:   AnthonyKenny98
% @Last Modified time: 2020-02-23 14:23:03

\todo[inline]{TODO: Provide information about this project's repository}

\chapter{Budget}
    \todo[inline]{Budget}

\chapter{RRT Supporting Documentation}
    \label{appendix:rrt_appendix}
    % @Author: AnthonyKenny98
% @Date:   2020-03-19 12:56:13
% @Last Modified by:   AnthonyKenny98
% @Last Modified time: 2020-04-06 10:20:38

\section{Justification of Modelling UAV as Prism}
\label{section:rrt_appendix_modelling}
    While it is possible for a \gls{UAV} to be modelled in precise detail, taking into account its exact shape, more often \glspl{UAV} are modelled as a 3D prism in motion planning problems, for the following reasons:
    \begin{itemize}
    \item It is a rare case that the negative space gained by modelling in such detail is utilised
    \item Representation of drone's configuration is much more complex.
    \item Computing edge collisions is much more computationally intensive.
    \end{itemize}

    % @Author: AnthonyKenny98
% @Date:   2020-04-05 12:55:29
% @Last Modified by:   AnthonyKenny98
% @Last Modified time: 2020-04-05 15:43:13

\begin{figure}[H]
\begin{centering}

\begin{tabular}{cc}
\begin{subfigure}{0.45\linewidth}
    \includegraphics[width=\linewidth]{appendices/rrt/img/DroneNegSpace.png}
    \caption{}
    \label{subfig:dronenegspace}
\end{subfigure}
\begin{subfigure}{0.45\linewidth}
    \includegraphics[width=\linewidth]{appendices/rrt/img/DroneRecPrism.png}
    \caption{}
    \label{subfig:dronerecprism}
\end{subfigure}
\end{tabular}
\caption[Modelling a UAV as a Rectangular Prism]{Modelling a \gls{UAV} as a Rectangular Prism. Red highlight demonstrates the model, overlayed over the exact schematic. Figure \ref{subfig:dronenegspace} shows how a drone can be modelled in high detail, but gains little useful free space when compared with Figure \ref{subfig:dronerecprism}, which models a drone as a rectangular prism.\cite{Thingbits}}
\label{fig:dronemodelling}
\end{centering}
\end{figure}

\section{Full Technical Specifications for RRT Implementation}
\label{section:rrt_appendix_tech_specs}
    % @Author: AnthonyKenny98
% @Date:   2020-04-05 14:09:51
% @Last Modified by:   AnthonyKenny98
% @Last Modified time: 2020-04-11 03:02:28

\begin{table}[H]
\begin{center}
\begin{tabular}{|p{.2\linewidth}|p{.74\linewidth}|}
    \hline
    \multicolumn{2}{|c|}{\textbf{General Specifications}} \\
    \hline
    \textbf{Requirement}             & \textbf{Description and Justification} \\
    \hline
    Implemented in C/C++    & 
        As outlined in Section \ref{subsection:project_structure}, the critical step in determining the design of specialized hardware to accelerate \gls{RRT} is CPU performance analysis of the algorithm to determine computational hot-spots. Implementations in C allow for the use of certain CPU profiling tools, unlike higher-level languages such as Python. \\
    \hline
    3D Workspace            & 
        The computational requirements of \gls{RRT} in \gls{3D} differs somewhat to that in \gls{2D}. Since autonomous \glspl{UAV} operate in 3D space, it was neccesary to have a \gls{3D} implementation to analyse. \\
    \hline
    \Gls{UAV} modelled in \gls{3D} as a rectangular prism  & 
        In theory, it is possible to model a \gls{UAV} much more precisely than a rectangular prism, taking into account its shape and negative space. However, in reality, modelling a \gls{UAV} as a \gls{3D} rectangular prism, defined by coordinates $\{x, y, z\}$ and Euler angles $\{\alpha, \beta, \gamma \}$, is more than sufficient (and more efficient). See Appendix \ref{section:rrt_appendix_modelling} for justification of this. \\
    \hline
    Mathematically Complete Collision Detection & 
        When \gls{RRT} is implemented for educational purposes, the edge collision calulations are often simplified to a sampling model which is \gls{probabilistically complete} but not \gls{mathematically complete}. In other words, it will catch most collisions by sampling a number of points along each edge, but there is always a possibility of an undetected collision. In real world applications, collisions must be calculated by method of geometric intersection to ensure all collisions are detected. \\
    \hline
\end{tabular}
\mycaption{General Technical Specifications for \gls{RRT} Implementation}{}
\label{table:RRT_Tech_Specs_General}
\end{center}
\end{table}

\begin{table}[H]
\begin{center}
\begin{tabular}{|p{.2\linewidth}|p{.74\linewidth}|}
    \hline
    \multicolumn{2}{|c|}{\textbf{Required Parameters}} \\
    \hline
    \textbf{Parameter}   & \textbf{Description and Justification} \\
    \hline
    $\epsilon$ (a.k.a. $\Delta q$) & 
        The maximum difference between two configurations. Larger values of $\epsilon$ can solve less obstacle dense problems faster, but take longer to solve problems with tight corners.\\
    \hline
    $K$ &
        The maximum number of configurations. This is largely correlative to the amount of time the user will allow the algorithm to run. Larger values of $K$ will take longer but generate better paths, while smaller values will execute for less time but generate more jagged paths or may not reach the goal node. The value of $K$ was varied to find the minimum execution time while still reaching the goal with high probability. \\
    \hline
    $DIM$ &
        The upper bound of each axis of a $DIM\times DIM\times DIM$ Workspace. Larger values leave more space to be explored, and thus require larger values of $K$ to reach the goal with high likelihood. \\
    \hline
    Goal Bias &
        The given probability that the graph will extend the graph $\epsilon$ distance from an existing configuration to a new configuration in the direction of the goal. \\
    \hline
\end{tabular}
\mycaption{Required Parameters for \gls{RRT} Implementation}{}
\label{table:RRT_Tech_Specs_Parameters}
\end{center}
\end{table}

\section{Assessment of Existing RRT Implementations}
\label{section:rrt_appendix_existing_implementations}
    % @Author: AnthonyKenny98
% @Date:   2020-04-05 13:32:41
% @Last Modified by:   AnthonyKenny98
% @Last Modified time: 2020-04-05 13:46:28

\begin{table}[H]
\begin{centering}
\begin{tabular}{|p{0.2\linewidth}|p{0.15\linewidth}|p{0.15\linewidth}|p{0.15\linewidth}|p{0.15\linewidth}|}
\hline
Repository  & Language  &  Workspace Dimension  & Object Model & Algorithmic Correctness \\
\hline
RoboJackets\cite{RoboJackets2019}   & \textcolor{mygreen}{C++}    &       &       & \\
\hline
\end{tabular}
\caption[Evaluation of Existing Open-Source Implementations of RRT]{Evaluation of Existing Open-Source Implementations of RRT. Links to Github repositories can be found in the Bibliography.}
\end{centering}
\end{table}

\section{Implementation of Key RRT Functions}
\label{section:rrt_appendix_function_impl}
    % @Author: AnthonyKenny98
% @Date:   2020-04-05 21:30:06
% @Last Modified by:   AnthonyKenny98
% @Last Modified time: 2020-04-11 03:04:30

\bigskip
\begin{algorithm}[H]
    \caption{\texttt{getRandomConfig()} as implemented for \gls{RRT}}
    \SetAlgoLined
    \SetArgSty{textnormal}
    \begin{tabular}{l l}
    \textbf{Inputs:}    & Dimensionality $N$,\\ 
                        & Upper Axis Bound $DIM$ \\
    \textbf{Output:}    & Random Configuration $q$ \\
    \end{tabular}

        $q.x \leftarrow$ randomFloat($DIM$) \\
        $q.y \leftarrow$ randomFloat($DIM$) \\
        $q.\alpha \leftarrow$ randomFloat($2\pi$) \\
        \If{$N == 3$}{
            $q.z \leftarrow$ randomFloat($DIM$) \\
            $q.\beta \leftarrow$ randomFloat($2\pi$) \\
            $q.\gamma \leftarrow$ randomFloat($2\pi$) \\
        }
        \Return $q$;\\
\end{algorithm}
\bigskip
Where \texttt{randomFloat(max)} returns a float between 0 and \texttt{max}.

\bigskip
\begin{algorithm}[H]
    \caption{\texttt{findNearestConfig()} as implemented for \gls{RRT}}
    \SetAlgoLined
    \SetArgSty{textnormal}
    \begin{tabular}{l l}
    \textbf{Inputs:}    & Graph $G$, \\
                        & New Configuration $q_{new}$ \\
    \textbf{Output:}    & Nearest Configuration $q_{nearest}$ \\
    \end{tabular}

        $q_{nearest} \leftarrow$ $G$.$q_{init}$ \\
        \For{$k = 0$ to $G$.existing\_nodes}{
            \If{distance($q_{new}$, $G.q$[$k$]) < distance($q_{new}$, $q_{nearest}$)}{
                $q_{nearest} \leftarrow$ $G.q$[$k$] \\
            }
        }
        \Return $q_{nearest}$ \\
\end{algorithm}
\bigskip
Where \texttt{distance($q_1$, $q_2$)} returns the Euclidean distance between two configurations.

\bigskip
\begin{algorithm}[H]
    \caption{\texttt{stepFromNearest()} as implemented for \gls{RRT}}
    \SetAlgoLined
    \SetArgSty{textnormal}
    \begin{tabular}{l l}
    \textbf{Inputs:}    & Configuration in Graph $q_{nearest}$,\\ 
                        & New Configuration $q_{new}$, \\
                        & Goal Bias $B$, \\
                        & Maximum Step Distance $\epsilon$, \\
                        & Graph $G$ \\
    \textbf{Output:}    & Updated New Configuration $q_{new}$ \\
    \end{tabular}

        \If{distance($q_{nearest}$, $q_{new}$) > $\epsilon$}{
            \If{randomFloat($1$) < $B$}{
                $q_{new} \leftarrow $ stepTowardConfig($q_{nearest}$, $G$.$q_{goal}$) \\
            } \Else{
                $q_{new} \leftarrow $ stepTowardConfig($q_{nearest}$, $q_{new}$) \\
            }
        }
        \Return $q_{new}$;\\
\end{algorithm}
\bigskip
Where \texttt{stepTowardConfig($q_1$, $q_2$)} returns a configuration $\epsilon$ from $q_1$ in the direction of $q_2$.

\bigskip
\begin{algorithm}[H]
    \caption{\texttt{configCollision()} as implemented for \gls{RRT}}
    \SetAlgoLined
    \SetArgSty{textnormal}
    \begin{tabular}{l l}
    \textbf{Inputs:}    & Dimensionality $N$,\\ 
                        & Occupancy Grid Map ($N$-Dimensional Array) $O$,\\ 
                        & Configuration $q$ \\
    \textbf{Output:}    & Boolean \\
    \end{tabular}

        \If{$N$ == 2 }{
            \Return $O$[gridLookup($q.x$)][gridLookup($q.y$)]
        }
        \Else{
            \Return $O$[gridLookup($q.x$)][gridLookup($q.y$)][gridLookup($q.z$)]
        }
\end{algorithm}
\bigskip
Where $O$ is a $N$-Dimensional array of booleans, with True representing an occupied grid and false representing an unoccupied one. \texttt{gridLookup()} is a function that maps a floating point coordinate to the correct integer of the grid in which it resides. For a map resolution of one, this is as simple as rounding a float down to an integer. \\

\begin{algorithm}[ht!]
    \caption{\texttt{configCollision()} as implemented for \gls{RRT} for 3D}
    \SetAlgoLined
    \SetArgSty{textnormal}
    \begin{tabular}{l l}
    \textbf{Inputs:}    & Edge $e$,\\ 
                        & Occupancy Grid Map (3-Dimensional Array) $O$,\\ 
                        & Maximum Step Distance $\epsilon$ \\
    \textbf{Output:}    & Boolean \\
    \end{tabular}
    
        $q_{min} \leftarrow $ minConfig($e.q_1$, $e.q_2$) \\
        \For{($x=  q_{min}.x$ to $q_{min}.x + \epsilon$)}{
            $q_{intersection} \leftarrow$ = edgeIntersectsPlane($e$, $x$) \\
            \If{$O$[$q_{intersection}.x$][$q_{intersection}.y$][$q_{intersection}.z$]}{
                \Return true
            }
        }
        \For{($y=  q_{min}.y$ to $q_{min}.y + \epsilon$)}{
            $q_{intersection} \leftarrow$ = edgeIntersectsPlane($e$, $y$) \\
            \If{$O$[$q_{intersection}.x$][$q_{intersection}.y$][$q_{intersection}.z$]}{
                \Return true
            }
        }
        \For{($z=  q_{min}.z$ to $q_{min}.z + \epsilon$)}{
            $q_{intersection} \leftarrow$ = edgeIntersectsPlane($e$, $z$) \\
            \If{$O$[$q_{intersection}.x$][$q_{intersection}.y$][$q_{intersection}.z$]}{
                \Return true
            }
        }
        \Return false
\end{algorithm}

While seemingly complex, the above algorithm merely steps through the mathematical process of checking the relevant $x$, $y$, and $z$ planes for a point of intersection with the edge $e$. It then looks up the \gls{OGM} $O$ to see if the grid corresponding with the point of intersection is occupied. If so, then it reports a collision by returning True. The function \texttt{edgeIntersectsPlane} follows the geometrical process of detecting a segment-plane intersection outlined in Appendix \ref{section:rrt_appendix_line_plane_intersection}. $q_{min}$ is calculated to be the origin point of the grid closest to the origin. In other words, the algorithm does not check for intersections throughout the entire map, only the maximum number of grids that could possible be intersected by the edge $e$, given the location of the two points of the edge, $e.p_1$ and $e.p_2$, and the maximum edge length $\epsilon$. The algorithm for \texttt{edgeCollision()} in \gls{2D} can be inferred from the above, checking segment-line intersections for $x$ and $y$ lines.

\newpage
\section{Geometrically Determining Segment-Plane Intersection}
\label{section:rrt_appendix_line_plane_intersection}
    % @Author: AnthonyKenny98
% @Date:   2020-04-06 09:45:02
% @Last Modified by:   AnthonyKenny98
% @Last Modified time: 2020-04-06 11:28:04
The method of edge collision detection in this project's implementation of \gls{RRT} relies on detecting segment-plane intersections. The planes are always set up to be parallel with either the $x-y$ plane, the $x-z$ plane, or the $y-z$ plane. Figure \ref{fig:parallelPlanes} demonstrates this point.

\begin{figure}[H]
\begin{centering}
\includegraphics[width=\linewidth]{appendices/rrt/img/planes}
\caption{Using Parallel Planes to determine Edge Collisions with Grids}
\label{fig:parallelPlanes}
\end{centering}
\end{figure}

A plane can be defined by 3 points, $P_a$, $P_b$, and $P_c$. In practice, the points defining a plane parallel to the x-y plane would have the following points:

    $$P_a = (x,y,z)$$
    $$P_b = (x+\Delta x,y,z)$$
    $$P_c = (x,y+\Delta y,z)$$

Two vectors, $\vec{AB}$ and $\vec{AC}$ can be determined.

The normal to the plane is the cross product: 
    $$\vec{AB} \times \vec{AC}$$

And the equation of the plane written as:

    $$a(x-x_0) + b(y-y_0) + c(z-z_0) = 0$$
    $$ax + by + cz = ax_0 + by_0 + cz_0$$

Where $<a,b,c>$ is the normal to the plane and $(x_0, y_0, z_0)$ is one of the points $P_a$, $P_b$, or $P_c$.
The RHS can be set to equal $d$, leaving:

    $$ax + by + cz = d$$

Now, the equation of a line can be written in the form:

    $$ax + by + cz = 0$$

And can be parameterized in the following form:
    $$\begin{cases}
        x = & x_1 + t(x_2 - x_1) \\
        y = & y_1 + t(y_2 - y_1) \\
        z = & z_1 + t(z_2 - z_1) \\

    \end{cases}$$

To find the point of intersection, we substitute the equation of the line into the equation of the plane, yielding:

    $$a(x_1 + t(x_2 - x_1)) + b(x_1 + t(x_2 - x_1)) + c(x_1 + t(x_2 - x_1)) = d$$

Rearranging to find an expression for $t$:

    $$t = \frac{d - (ax_1 + by_1 + cz_1)}{a(x_2-x_1) + b(y_2-y_1) + c(z_2-z_1)}$$

Knowing $t$, we can find the point of intersection, $P_X$ to be:

    $$\begin{cases}
        x_X(t) = & x_1 + t(x_2 - x_1) \\
        y_X(t) = & y_1 + t(y_2 - y_1) \\
        z_X(t) = & z_1 + t(z_2 - z_1) \\
    \end{cases}$$

Finally, the following equalities are evaluated to see if the point lies on the segment:

    $$x_1 \leq x_X \leq x_2$$
    $$y_1 \leq y_X \leq y_2$$
    $$z_1 \leq z_X \leq z_2$$

If so, then the grids corresponding to the point of intersection can be marked as intersected.

\section{Timing Methodology of RRT Analysis}
\label{section:rrt_appendix_timing}
    % @Author: AnthonyKenny98
% @Date:   2020-04-06 08:00:17
% @Last Modified by:   AnthonyKenny98
% @Last Modified time: 2020-04-06 15:28:48

\subsubsection*{VTune Profiler}
    VTune Profiler is an application for software performance analysis. It provides functionality to examine hot-spots for CPU execution time through a top down analysis. The top down analysis tool shows the percentage of CPU time taken up by each function. It was used to initially profile the algorithm's performance.

    \subsubsection{Internal Timing}
        There are several limitations to using VTune Profiler. First, it can only profile software running on Intel processors, which implement the x86 \gls{ISA}. In anticipation of potentially needing to run performance analysis on a \gls{RISC-V} processor, another method was required. Secondly, VTune Profiler takes a long time to run, as it needs to conduct a lot of analysis that is extraneous to the purpose of this thesis. This became prohibitive when it came to conducting hundreds of tests for different parameterizations, with each test running \gls{RRT} a minimum of 100 times. Finally, it was not customizable to ignore certain parts of the implementation, such as logging functionality. While the implementation was designed in such a way that these should not intefere, it led to a lot of irrelevant data. A simple and effective alternative for measuring execution performance was to insert timing functionality into the software itself.

        Internal timing was implemented based on the inbuilt C \texttt{clock()} function and \\ \texttt{CLOCKS\_PER\_CYCLE} macro, and wrapping each function of interest in a performance tracking struct. This can be seen in the project's \gls{RRT} sub-repository under \texttt{performance.h}.

    \subsubsection*{Comparison}
        Before proceeding to use the internal timing method, it was important to verify that this method yielded similar results to VTune Profiler for the same program. Table \ref{table:timing_calibration} summarizes the results of analysis of a simple C executable. The program calls 5 functions, $\{A, B, C, D, E\}$, each a simple iteration in which a integer is incremented. Since the Internal Timing method returned similar results to the (trusted) VTune Profiler, it was considered to be a reliable method. While it was encouraging to see both methods returned similar results for absolute execution time, the more important metric was the similarity in percentage of total execution time. For good measure, a $\chi^2$ test of hypothesis was conducted and for one degree of freedom showed more that acceptable results.

        \begin{table}[H]
        \begin{center}
        \begin{tabular}{|m{0.09\linewidth}|m{0.16\linewidth}|m{0.16\linewidth}|m{0.16\linewidth}|m{0.16\linewidth}|m{0.09\linewidth}|}
        \hline
        \multirow{2}{*}{function} & \multicolumn{2}{c|}{Vtune Profiler} & \multicolumn{2}{c|}{Internal Timing}  & \multirow{2}{*}{$\chi^2$}\\
        \cline{2-5}
                    & time (s)  & time (\% total) & time (s) & time (\% total)  &   \\
        \hline
        A       & 0.488     & 57.4\%        & 0.497 & 57.6\%  &  0.00016       \\
        B       & 0.2       & 23.5\%        & 0.198 & 23.1\%  &  0.00002       \\
        C       & 0.102     & 12.0\%        & 0.099 & 11.5\%  &  0.00009       \\
        D       & 0.048     & 5.7\%         & 0.049 & 5.6\%   &  0.00002       \\
        E       & 0.012     & 1.4\%         & 0.019 & 2.2\%   &  0.00408       \\
        \hline
        \end{tabular}
        \end{center}
        \caption{Comparison of Timing Methods}
        \label{table:timing_calibration}
        \end{table}
\newpage
\section{Justification of K:DIM Ratio}


\end{appendices} 
\clearpage 
\listoftodos 

\end{document}
